% =====================================================================
% figures_snippet.tex  --  paste these 4 figure blocks into your paper.
%
% Assumes: images are uploaded to a folder "figures/" that sits in the SAME
% directory as main.tex, with the exact filenames below, and that the preamble
% of main.tex has:  \usepackage{graphicx}
%
% Figure numbers are AUTOMATIC: LaTeX numbers them in order of appearance, so
% they slot in correctly alongside the other figures already in your paper.
% Refer to them in text with \ref{fig:...} (labels are set below).
% Each figure is followed by a short explanation paragraph you can keep or move.
% =====================================================================

\begin{figure}[t]
  \centering
  \includegraphics[width=\linewidth]{figures/LLaVA_top10.png}
  \caption{Per-token contrastive adjustment $d = E - A$ over the expert's
  top-10 object candidates, for LLaVA-1.5-7B. Left column: real objects; right
  column: hallucinated objects. Top row: VCD; bottom row: SID. Red bars are
  amplified ($d > 0$), blue bars are suppressed ($d < 0$); the dashed line marks
  the group mean. The adjustment is small and positive for both groups, so
  neither is suppressed.}
  \label{fig:d-llava-top10}
\end{figure}

\paragraph{Figure~\ref{fig:d-llava-top10} (LLaVA, top-10).}
On LLaVA the contrastive adjustment $d = E - A$ is small and \emph{positive} for
both groups (VCD: $+0.133$ real, $+0.164$ hallucinated; SID: $+0.218$ real,
$+0.114$ hallucinated), so contrastive decoding on average \emph{amplifies}
object logits rather than suppressing them. Hallucinated objects are not pushed
down relative to real ones: under VCD they are amplified slightly more, and under
SID both are amplified with no meaningful gap. We infer that on LLaVA the
contrastive term gives no selective suppression of hallucinations; its effect on
object tokens is a small, non-discriminative upward nudge.

\begin{figure}[t]
  \centering
  \includegraphics[width=\linewidth]{figures/Qwen_top10.png}
  \caption{Per-token contrastive adjustment $d = E - A$ over the expert's
  top-10 object candidates, for Qwen2.5-VL-7B. Layout as in
  Figure~\ref{fig:d-llava-top10}. Real objects are suppressed more than
  hallucinations (the real-object panels are more negative), the opposite of the
  intended behaviour.}
  \label{fig:d-qwen-top10}
\end{figure}

\paragraph{Figure~\ref{fig:d-qwen-top10} (Qwen, top-10).}
On Qwen the adjustment is \emph{negative} overall, so contrastive decoding lowers
object logits, but it lowers \emph{real} objects far more than hallucinations
(VCD: $-0.984$ real vs.\ $-0.429$ hallucinated; SID: $-0.674$ vs.\ $-0.172$), and
a larger share of hallucinated tokens is amplified ($36\%$ vs.\ $25\%$ positive
under VCD). We infer that the suppression is aimed the wrong way: the tokens it
pushes down hardest are the correct ones, while hallucinations are comparatively
spared, which is the opposite of the method's stated goal.

\begin{figure}[t]
  \centering
  \includegraphics[width=\linewidth]{figures/LLaVA_emitted.png}
  \caption{Contrastive adjustment $d = E - A$ on emitted object words (exact
  caption-level labels; one thin spike per emitted word), for LLaVA-1.5-7B.
  Layout as in Figure~\ref{fig:d-llava-top10}. Hallucinated words are amplified
  more than real words.}
  \label{fig:d-llava-emitted}
\end{figure}

\paragraph{Figure~\ref{fig:d-llava-emitted} (LLaVA, emitted words).}
Restricting to the object words the model actually wrote, where the labels are
exact, sharpens the picture: emitted hallucinated words receive a much larger
positive adjustment than real words (VCD: $+0.789$ vs.\ $+0.264$; SID: $+0.684$
vs.\ $+0.296$), roughly a threefold stronger amplification of hallucinations. We
infer that, far from suppressing the hallucinations it emits, contrastive
decoding boosts them more than it boosts the real objects.

\begin{figure}[t]
  \centering
  \includegraphics[width=\linewidth]{figures/Qwen_emitted.png}
  \caption{Contrastive adjustment $d = E - A$ on emitted object words (exact
  caption-level labels; one thin spike per emitted word), for Qwen2.5-VL-7B.
  Layout as in Figure~\ref{fig:d-llava-top10}. Real objects are strongly
  suppressed while hallucinations sit near zero or are amplified.}
  \label{fig:d-qwen-emitted}
\end{figure}

\paragraph{Figure~\ref{fig:d-qwen-emitted} (Qwen, emitted words).}
On Qwen the emitted-word view is the starkest: real object words are strongly
suppressed while hallucinated words sit near zero or are amplified (VCD:
$-1.281$ real vs.\ $-0.097$ hallucinated; SID: $-0.747$ real vs.\ $+0.429$
hallucinated, where hallucinations are pushed \emph{up} as real objects are
pushed down). We infer that the contrastive term actively works against
grounding here, penalising the correct objects the model emits while leaving or
making hallucinations more likely, consistent with Qwen's degraded CHAIR scores.


% =====================================================================
% Proxy A methodology + Gaussian-noise parameter table.
% Table number is automatic; refer to it with \ref{tab:proxyA-noise}.
% =====================================================================

\paragraph{Proxy A: matched Gaussian noise.} Proxy~A tests whether the effect of
contrastive decoding can be reproduced by adding random noise of a matched size
to the expert model alone. We discard the amateur branch entirely and run a
single expert-only greedy decode. At each decoding step we take the plain expert
logits $E$ and add, to the logits of the object-category tokens, independent
Gaussian noise drawn from $\mathcal{N}(\mu,\sigma^{2})$. The mean $\mu$ and
standard deviation $\sigma$ are set to the measured mean and standard deviation
of the contrastive adjustment $d = E - A$ for the corresponding method
(Table~\ref{tab:proxyA-noise}), so the injected noise has the same magnitude the
real method adds. We then apply the same adaptive plausibility constraint
(threshold $\log\beta + \max_{w} E_{t}(w)$, with $\beta = 0.2$, computed from the
expert logits) and take the greedy token. The proxy therefore never sees a
degraded image and performs no contrast; its only ingredients are the expert
logits and matched Gaussian noise. We score the resulting captions on the
standard CHAIR benchmark and compare them against the real VCD and SID runs.

\begin{table}[t]
  \centering
  \small
  \setlength{\tabcolsep}{6pt}
  \caption{Proxy~A against the real methods on CHAIR. Greedy is the undefended
  baseline; VCD and SID are the real contrastive methods. Proxy~A discards the
  amateur branch and instead injects Gaussian noise $\mathcal{N}(\mu,\sigma^{2})$
  into the expert's object-category logits, with $(\mu,\sigma)$ (shown in
  brackets) matched to the measured adjustment $d = E - A$ of VCD. CHAIR-S and
  CHAIR-I are percentages, lower is better ($\downarrow$). Proxy~A reproduces the
  CHAIR level of the real methods without any contrast.}
  \label{tab:proxyA-noise}
  \begin{tabular}{ll cc}
    \toprule
    Model & Setting & CHAIR-S~$\downarrow$ & CHAIR-I~$\downarrow$ \\
    \midrule
    LLaVA-1.5-7B  & Greedy    & $50.0$ & $13.9$ \\
                  & VCD (Real) & $55.0$ & $15.6$ \\
                  & SID (Real) & $54.2$ & $15.1$ \\
                  & Proxy A ($\mu{=}{+}0.199,\ \sigma{=}0.853$) & $52.4$ & $14.8$ \\
    \midrule
    Qwen2.5-VL-7B & Greedy    & $30.8$ & $7.7$ \\
                  & VCD (Real) & $36.4$ & $10.5$ \\
                  & SID (Real) & $38.4$ & $10.5$ \\
                  & Proxy A ($\mu{=}{-}0.846,\ \sigma{=}1.531$) & $30.8$ & $8.2$ \\
    \bottomrule
  \end{tabular}
\end{table}

\paragraph{What Table~\ref{tab:proxyA-noise} shows.} Neither real method reduces
hallucination below greedy: on LLaVA both VCD and SID sit at or above the greedy
baseline, and on Qwen they raise CHAIR clearly above it. Proxy~A, which uses no
amateur branch and no contrast but only matched Gaussian noise, lands in the same
CHAIR band as the real methods on LLaVA ($52.4$ vs.\ $55.0$ and $54.2$) and stays
near greedy on Qwen. We infer that on LLaVA the effect of contrastive decoding is
indistinguishable from adding random noise of the same magnitude, so the amateur
branch contributes no targeted hallucination reduction, while on Qwen the real
methods are in fact worse than matched noise. In neither case does contrastive
decoding provide a benefit that a simple noise injection cannot account for.
